\begin{claim}{F.29}{Proved}{All finite affine marginal images on the four-parameter product-one locus}
Retain the actual bases $Y_i^-=X_i$ and
$Y_i^+=\mathcal B_+(X_i)$, where $\mathcal B_+$
is the based isomorphism of F.23.13.
For every finite field $k$ and continuous
$\rho:V_\pm\to\operatorname{Aff}_1(k)$ with
scalar character $\chi$, the exact locus
\[
 \chi(Y_5^\pm)=\chi(Y_6^\pm)=1,
 \qquad \prod_{i=1}^7\chi(Y_i^\pm)=1
                                                               \tag{F.29.1}
\]
has the image equality
\[
                         \rho(\mathcal K_\pm)=\rho(R_\pm).
                                                               \tag{F.29.2}
\]
Every characteristic, finite scalar-field extension
degree and larger original coefficient field is included.

Write a negative character in this locus as
$(r,s,t,u,1,1,c^{-1})$, where $c=rstu$.
On its nontrivial-scalar part
$C=\langle u,c^{-1}\rangle\ne\{1\}$, put
$k_1=\mathbb F_p(u,c)$, $k_0=\mathbb F_p(r,s,t,u)$,
$q_1=|k_1|$, $m=[k_0:k_1]$ and
$\chi_j=\chi^{q_1^j}$, $0\leq j<m$.
The summed continuous $H^1$ restriction to
$K=\mathcal K_-$ is injective. At the cocycle
level its kernel and rank are exactly
\[
 \begin{gathered}
 \left\{\bigl(b_j(\chi_j-1)\bigr)_{j=0}^{m-1}:
                 b_j\in k_0,\ \sum_jb_j=0\right\},\\
                       \operatorname{rank}=6m+1
                 \quad\hbox{on }7m\hbox{ coordinates}.
 \end{gathered}                                                \tag{F.29.3}
\]
\end{claim}

\begin{proof}
\emph{Actual fields, completed cases and the four remaining charts.}
Work first with $K=\mathcal K_-\leq R=R_-\leq V=V_-$.
F.21 and the actual marginal surjection of F.22
give $\overline{\langle\!\langle K\rangle\!\rangle}_R=R$.
Apply this equality to the continuous abelian
map $\chi|R$. Normal generation in its abelian
image gives $\chi(K)=\chi(R)$.
Since the actual deleted generators are
$X_4,X_5,X_6,X_7$, their scalars now give
\[
 \begin{gathered}
 C=\chi(K)=\chi(R)=\langle u,c^{-1}\rangle,
 \qquad D=\chi(V)=\langle r,s,t,u\rangle,\\
 k_1=\mathbb F_p(C)=\mathbb F_p(u,c),\qquad
 k_0=\mathbb F_p(D)=k_1(r,t),\qquad s=c/(rtu).
 \end{gathered}                                                \tag{F.29.4}
\]
All displayed character values are nonzero.
For $C=\{1\}$ the image of $R$ is a translation
group, hence abelian. The same normal-generation
map gives (F.29.2) directly. This case is precisely
$u=c=1$.
For $C\ne\{1\}$ and $k_0=k_1$, F.25 gives the
affine equality, while F.24 gives (F.29.3) with
$m=1$.

For $u=1$, the original character is exactly
$(r,s,t,1,1,1,(rst)^{-1})$ in the unchanged basis.
F.28 supplies the affine equality and, on its
nontrivial-scalar part, the stated restriction
kernel. The coefficient homomorphism making this
comparison is $u\mapsto1$, fixing $r,s,t$ and
sending $c$ to $rst$. It sends each original
word evaluation to that of F.28, and sends the
principal-pivot denominator $1-c$ to $1-rst$.
Thus the comparison uses the actual source maps.

We now work with $u\ne1$, $C\ne\{1\}$ and
$k_1\subsetneq k_0$. The four disjoint coordinate
charts are $c\ne1,r\ne1$; $c\ne1,r=1$;
$c=1,r\ne1$; and $c=1,r=1$.
We give the actual principal quotient, invariant
maps and source calculations in each chart.

\emph{Coefficient maps, principal quotients and Frobenius.}
The first two charts use the coefficient homomorphism
\[
 \begin{aligned}
 \mathbb Z[t_1^{\pm1},\ldots,t_7^{\pm1}]&\longrightarrow
     A_4=\mathbb Z[r^{\pm1},s^{\pm1},t^{\pm1},u^{\pm1}],\\
 (t_1,t_2,t_3,t_4,t_5,t_6,t_7)&\longmapsto
                  (r,s,t,u,1,1,(rstu)^{-1}).
 \end{aligned}                                                \tag{F.29.5}
\]
An exponent vector $(a_i)$ maps to
$r^{a_1-a_7}s^{a_2-a_7}t^{a_3-a_7}u^{a_4-a_7}$.
The chart $r=1$ further applies the homomorphism
fixing $s,t,u$ and sending $r$ to $1$.
All these are derivation identities on the
specified discrete free words. Evaluating in
a finite field gives the actual cocycle rule
$d(ab)=d(a)+\chi(a)d(b)$.
Every seven-tuple of generator values defines
a homomorphism to the finite affine group and
extends uniquely to $V=\widehat F_7$.
Thus the continuous cocycle space is $k_0^7$;
the Laurent calculations are not asserted to
define continuous cocycles with infinite coefficient ring.

For the character in a given chart put
$\delta=(\chi(X_i)-1)_{i=1}^7$.
Use pivot $p=7$ when $c\ne1$, and pivot $p=4$
when $c=1,u\ne1$. Let $I_p$ be the increasing
list of the other six indices. Define
\[
 \begin{aligned}
 P_p(v)&=(v_i-\delta_i\delta_p^{-1}v_p)_{i\in I_p},
       &S_p(w)_p&=0,\quad S_p(w)_{I_p}=w,\\
 P_pS_p&=I_6,
       &S_pP_p(v)&=v-\delta\delta_p^{-1}v_p.
 \end{aligned}                                                \tag{F.29.6}
\]
Here $\delta_7=(1-c)/c\ne0$ in the first charts;
the last charts have $u=h^{-1}$, $h=rst\ne1$,
and $\delta_4=(1-h)/h\ne0$.
These two composites identify the actual continuous
cohomology quotient $Q_\chi=H^1_{\rm cont}(V,k_{0,\chi})$
with $k_0^6$ in each displayed coordinate order.

Retain the eleven actual F.23 automorphisms
$B_l(v)=u_l\alpha_l(v)u_l^{-1}$, $0\leq l\leq10$.
They and their inverses preserve $K$ and fix
the character. Their cocycle matrices $A_l$
have rows $D(B_l(X_i))$, so $A_l\delta=\delta$.
Put $\overline A_l=P_pA_lS_p$. For each basis
$J$ below the cohomology matrix
$T_l=J^{-1}\overline A_lJ$ is certified by the
integral identity
\[
 \delta_p(A_lS_pJ-S_pJT_l)
             =\delta\,\bigl(\operatorname{row}_p(A_lS_pJ)\bigr).
                                                               \tag{F.29.7}
\]
This clears precisely the displayed principal
pivot and proves the claimed action after every
permitted field evaluation.
All source indices in this proof begin at zero:
$0=\texttt{old0}$, $1,\ldots,5$ are
$\texttt{old\_1\_0},\ldots,\texttt{old\_1\_4}$,
and $6,\ldots,10$ are
$\texttt{old\_-1\_0},\ldots,\texttt{old\_-1\_4}$.
Their positions in an enumeration beginning at
one are $l+1$; the word labels remain unchanged.

For $P_0=X_1\cdots X_7$, the product-one condition
makes $d(P_0)$ a covector on cohomology.
The pure Artin map fixes $P_0$, and direct
conjugation gives $d(B_l(P_0))=\chi(u_l)d(P_0)$.
Also each character-one meridian has
$d(B_l(X_i))=\chi(u_la_{li})d(X_i)$, where
$\alpha_l(X_i)=a_{li}X_ia_{li}^{-1}$ is its
actual pure braid conjugator.
These boundary and meridian evaluations are
the quotient covectors used below.

Put $r_j=r^{q_1^j}$, $s_j=s^{q_1^j}$,
$t_j=t^{q_1^j}$; $u,c\in k_1$ stay fixed.
The exact coefficient map
\[
 d\longmapsto d^{q_1^j},\qquad
 k_{0,\chi}\longrightarrow k_{0,\chi_j},\qquad
 Q_\chi\longrightarrow Q_j=H^1_{\rm cont}(V,k_{0,\chi_j})
                                                               \tag{F.29.8}
\]
is additive and semilinear and preserves cocycles,
principals and continuity. Its inverse is the
power $q_1^{m-j}$, with the identity at $j=0$.
It commutes with every actual source precomposition.
Thus each displayed basis, inclusion, quotient,
source identity and inverse map transports by
the indicated Frobenius substitutions.
All $\chi_j$ agree on $K$, by (F.29.4).
Moreover
\[
                     (r_j,t_j)=(r_l,t_l)\quad\Longrightarrow\quad j=l.
                                                               \tag{F.29.9}
\]
Indeed such an equality for distinct indices
would make a nonidentity $k_1$-Frobenius power
fix the generators of $k_0=k_1(r,t)$.
For $0<a<m$, the polynomial $X^{q_1^a}-X$ has
degree smaller than $|k_0|$, and cannot vanish
on the whole field. In either $r=1$ chart,
$k_0=k_1(t)$, so the $t_j$ are distinct and
lie outside $k_1$. In particular $t_j^2\ne1$.
In either $r\ne1$ chart every $r_j\ne1$.

\emph{Chart $c\ne1$, $r\ne1$: the three-dimensional core.}
Use $p=7$ and the six coordinates of
$X_1,\ldots,X_6$. Put $b=(1,r,rs,rst,c,c)$.
The boundary and the actual $X_5,X_6$ evaluations
give the explicit maps
\[
 \begin{aligned}
 \pi(v)&=(bv,v_5,v_6),\\
 \iota(x,y,z)&=(-rx-rsy-rstz,x,y,z,0,0),\\
 \sigma(z_\partial,z_5,z_6)
             &=(z_\partial-c(z_5+z_6),0,0,0,z_5,z_6).
 \end{aligned}                                                \tag{F.29.10}
\]
They satisfy $\pi\sigma=I_3$,
$\ker\pi=\operatorname{im}\iota=:W$ and
$\iota(v_2,v_3,v_4)+\sigma\pi(v)=v$.
In particular the basis and its inverse are
\[
 \begin{gathered}
 J=(\iota e_1,\iota e_2,\iota e_3,e_1,e_5-ce_1,e_6-ce_1),\\
 J^{-1}(v)=(v_2,v_3,v_4,bv,v_5,v_6).
 \end{gathered}                                                \tag{F.29.11}
\]
Here the $e_i$ inside $\iota$ are the three
domain coordinate vectors; the remaining ones
are the six quotient coordinate vectors.
The displayed decomposition proves both basis
composites directly.
In this basis every actual source matrix is
$\left(\begin{smallmatrix}G_l&E_l\\0&D_l\end{smallmatrix}\right)$,
where $G_l$ acts on $W$ and $D_l$ is diagonal
on the boundary, $X_5$ and $X_6$ coordinates.
Equivalently $\overline A_l\iota=\iota G_l$ and
$\pi\overline A_l=D_l\pi$.
The actual source-word table is
\[
 \begin{array}{c|c|ccc}
 l&G_l&\lambda_{\partial,l}&\lambda_{5,l}&\lambda_{6,l}\\\hline
 0&rI_3&1&r&r\\
 2&*&rt&rt/u&rct\\
 7&tI_3&t^{-1}&t/u&ct\\
 8&*&1&1&t^{-1}\\
 9&*&1&c&1
 \end{array}
                                                               \tag{F.29.12}
\]
The stars denote retained full core matrices;
the additional exact identities used for them are
\[
 (G_8-I_3)^2(G_8-(c/u)I_3)=0,
 \qquad (G_9-I_3)^2(G_9-(tu)^{-1}I_3)=0.
                                                               \tag{F.29.13}
\]
These matrix products are integral Laurent
identities verified from the literal source words.

We explain the invariant-subspace classification
used in this and the next two separating charts.
Take a nonzero source-invariant subspace
$U\leq Q_j$ with no proper nonzero invariant
subspace. Its intersection with the displayed
$W_j$ is invariant. A nonzero intersection
therefore gives $U\leq W_j$.
For zero intersection, $\pi_j|U$ is injective.
A nonzero coordinate projection of this map
has invariant kernel, so is injective as well.
It identifies $U$ with one of the displayed
one-dimensional quotient characters. This
uses the actual coordinate projections and
does not require distinct quotient characters
inside a single factor.

For $j\ne l$, compare these subspaces through
any source-intertwining isomorphism. Source
operators acting by scalars $a,b$ on its two
sides give $(a-b)F=0$, so $F\ne0$ forces $a=b$.
In the present chart source $0$ separates a
boundary type from every other type, because
its scalars are $1$ and $r_l\ne1$.
Two core subspaces force $r_j=r_l,t_j=t_l$
through sources $0,7$, contrary to (F.29.9).
Two boundary types force $t_j=t_l$ through
source $7$ and then $r_j=r_l$ through source $2$.
For the same type $X_5$ or $X_6$, sources $0,7$
determine the pair $(r_j,t_j)$ after cancelling
the common units $u$ or $c$.
The different types $X_5,X_6$ are separated
by source $9$, with scalars $c\ne1$ and $1$.

A core subspace in factor $j$ isomorphic to
type $X_5$ in factor $l$ must be one-dimensional.
Source $9$ and (F.29.13) force
$(c-1)^2(c-(t_ju)^{-1})=0$, hence
$t_j=(cu)^{-1}\in k_1$.
Inverse Frobenius makes $t=(cu)^{-1}$ and all
$t_l=t_j$. Source $0$ then forces $r_j=r_l$,
contradicting (F.29.9).
For a core subspace versus type $X_6$, source
$8$ and (F.29.13) force
$t_l^{-1}\in\{1,c/u\}$, so
$t_l\in\{1,u/c\}\subset k_1$.
Again all $t_j=t_l$, and source $0$ contradicts
(F.29.9). These comparisons cover every pair
of the four types, including the same-type
pairs. The inverse of an intertwining
isomorphism covers the opposite orientation.

\emph{Chart $c\ne1$, $r=1$: two core coordinates and four quotients.}
Here $c=stu$, $k_1=\mathbb F_p(u,c)$ and
$k_0=k_1(t)$. In the same principal quotient
$p=7$, let $b=(1,1,s,st,c,c)$ and use
\[
 \begin{gathered}
 \pi(v)=(bv,v_1,v_5,v_6),\qquad
 \iota(x,y)=(0,-sx-sty,x,y,0,0),\\
 \sigma(z_\partial,z_1,z_5,z_6)
       =(z_1,z_\partial-z_1-c(z_5+z_6),0,0,z_5,z_6),\\
 J=(\iota e_1,\iota e_2,e_2,e_1-e_2,e_5-ce_2,e_6-ce_2),\\
 J^{-1}(v)=(v_3,v_4,bv,v_1,v_5,v_6).
 \end{gathered}                                                \tag{F.29.14}
\]
The identities $\pi\sigma=I_4$,
$\ker\pi=\operatorname{im}\iota=W$ and
$\iota(v_3,v_4)+\sigma\pi(v)=v$ prove the
maps and both basis inverses explicitly.
All eleven source matrices again have the
displayed invariant-core and diagonal-quotient
form. Their required entries are
\[
 \begin{array}{c|c|cccc}
 l&G_l&\lambda_{\partial,l}&\lambda_{1,l}
                         &\lambda_{5,l}&\lambda_{6,l}\\\hline
 1&*&1&u^{-2}&1&1\\
 2&tI_2&t&t^{-1}&t/u&ct\\
 5&I_2&1&c^2&t^{-1}&u^{-1}\\
 7&tI_2&t^{-1}&t&t/u&ct\\
 8&*&1&1&1&t^{-1}\\
 9&*&1&1&c&1
 \end{array}
                                                               \tag{F.29.15}
\]
The smaller core satisfies the exact quadratic
identities
\[
 (G_8-I_2)(G_8-(c/u)I_2)=0,
 \qquad (G_9-I_2)(G_9-(tu)^{-1}I_2)=0.
                                                               \tag{F.29.16}
\]
They, the basis and every quotient coordinate
are independently checked from all eleven
literal source words after $r\mapsto1$.

The five types are now $W,\partial,X_1,X_5,X_6$.
Their source-$2$, source-$7$ scalar pairs are,
respectively,
$(t_j,t_j)$, $(t_j,t_j^{-1})$,
$(t_j^{-1},t_j)$, $(t_j/u,t_j/u)$,
and $(ct_j,ct_j)$.
The same-type comparisons determine $t_j=t_l$,
which is impossible for $j\ne l$.
A core type versus boundary or $X_1$ would
force $t_l^2=1$ through these two sources.
A boundary or $X_1$ type versus $X_5$ or $X_6$
would likewise force $t_j^2=1$ on the first
side. These conclusions contradict
$t_j\notin k_1$.
Types $X_5,X_6$ are separated by source $9$.
A core type versus $X_5$ would force
$c=(t_ju)^{-1}$ by (F.29.16) and source $9$,
putting $t_j\in k_1$.
A core type versus $X_6$ would force
$t_l\in\{1,u/c\}\subset k_1$ by source $8$.
Both are impossible.

The only pair left is boundary in factor $j$
versus $X_1$ in factor $l$. Sources $1,5$ force
$u^2=c^2=1$. Since $c\ne1$ and $u\ne1$,
the factorization $(a-1)(a+1)=a^2-1$ in the
field gives
\[
 c=u=-1,\qquad p\ne2,\qquad
 s=t^{-1},\qquad k_1=\mathbb F_p,\qquad
 t\notin\mathbb F_p.
                                                               \tag{F.29.17}
\]
In characteristic two $a^2=1$ has only the
root $1$, so this remaining case is absent.

In (F.29.17) the coordinate order of (F.29.14)
is $(x,y,z_\partial,z_1,z_5,z_6)$.
Sources $2,7$ have core $tI_2$ and quotient
diagonals $(t,t^{-1},-t,-t)$ and
$(t^{-1},t,-t,-t)$, respectively.
The nonzero factors $t-t^{-1}$ and $2t$ follow
from $t\notin\mathbb F_p$ and $p\ne2$.
A boundary-character eigenvector must therefore
have its $X_1,X_5,X_6$ coordinates zero, by
source $2$. Its boundary coordinate cannot be
zero, because then source $7$ would compare
$t$ on the core with $t^{-1}$ on that vector.
Normalize this nonzero coordinate to one.
Source $7$ then uniquely forces the first
vector in the following exact identities:
\[
 \begin{gathered}
 v_\partial=(1-2t,2,1,0,0,0)^{\mathsf T},
       \qquad T_7v_\partial=t^{-1}v_\partial,\\
 v_1=(-t,0,0,1,0,0)^{\mathsf T},
       \qquad T_2v_1=t^{-1}v_1,\\
 (T_1-I_6)v_\partial=(T_1-I_6)v_1
                         =(3t-1,-2,0,0,0,0)^{\mathsf T}.
 \end{gathered}                                                \tag{F.29.18}
\]
The boundary scalar at source $1$ is $1$.
Its required eigenvector equation contradicts
the nonzero second coordinate $-2$ in
(F.29.18). Thus the boundary-character line
does not occur as a submodule in this case.
For completeness, source $7$ first kills the
other quotient coordinates of an $X_1$ line;
source $2$ uniquely forces $v_1$ after its
$X_1$ coordinate is normalized. The same
residual excludes that line too.
The coefficient specialization for these
identities is $r\mapsto1$, $s\mapsto t^{-1}$,
$u\mapsto-1$. A monomial $s^a t^b u^d$ maps
to $(-1)^d t^{b-a}$, and the pivot becomes
$\delta_7=-2$. The separate literal three-word
audit verifies (F.29.7) with this exact pivot
for sources $1,2,7$, before using its proved
odd-characteristic invertibility.
The same argument transports by Frobenius
to every factor. It excludes the last possible
cross-type comparison and completes this chart.

\emph{Chart $c=1$, $r\ne1$: the actual pivot $X_4$.}
Put $h=rst=u^{-1}\ne1$.
The character is now literally
$(r,s,t,h^{-1},1,1,1)$.
Its coefficient map is
\[
 \begin{aligned}
 \mathbb Z[t_1^{\pm1},\ldots,t_7^{\pm1}]&\longrightarrow
       A_3=\mathbb Z[r^{\pm1},s^{\pm1},t^{\pm1}],\\
 (t_1,t_2,t_3,t_4,t_5,t_6,t_7)&\longmapsto
                         (r,s,t,(rst)^{-1},1,1,1).
 \end{aligned}                                                \tag{F.29.19}
\]
The exponent image is
$r^{a_1-a_4}s^{a_2-a_4}t^{a_3-a_4}$.
This is also the homomorphism
$A_4/(rstu-1)\to A_3$ sending $u$ to
$(rst)^{-1}$ and fixing $r,s,t$.
The reverse map sends $r,s,t$ to their classes;
the relation $rstu=1$ makes the two composites
the identity, including on $u$.
The maps are applied to the unlocalized word
identities, followed here by localization at
$1-h$ for the actual pivot $X_4$ in (F.29.6).
The scalar fields are $k_1=\mathbb F_p(h)$
and $k_0=k_1(r,t)$ because $s=h/(rt)$.

The six quotient coordinates are now those
of $X_1,X_2,X_3,X_5,X_6,X_7$, in that order.
Put $b=(1,r,rs,1,1,1)$. The exact maps are
\[
 \begin{aligned}
 \pi(v)&=(bv,v_4,v_5,v_6),\\
 \iota(x,y)&=(-rx-rsy,x,y,0,0,0),\\
 \sigma(z_\partial,z_5,z_6,z_7)
              &=(z_\partial-z_5-z_6-z_7,0,0,z_5,z_6,z_7).
 \end{aligned}                                                \tag{F.29.20}
\]
They satisfy $\pi\sigma=I_4$,
$\ker\pi=\operatorname{im}\iota=W$ and
$\iota(v_2,v_3)+\sigma\pi(v)=v$.
Thus the basis maps are explicitly
\[
 \begin{gathered}
 J=(\iota e_1,\iota e_2,e_1,e_4-e_1,e_5-e_1,e_6-e_1),\\
                    J^{-1}(v)=(v_2,v_3,bv,v_4,v_5,v_6).
 \end{gathered}                                                \tag{F.29.21}
\]
Every actual source has invariant two-dimensional
core and diagonal quotient coordinates
$\partial,X_5,X_6,X_7$. The literal table is
\[
 \begin{array}{c|c|cccc}
 l&G_l&\lambda_{\partial,l}&\lambda_{5,l}
                         &\lambda_{6,l}&\lambda_{7,l}\\\hline
 0&rI_2&1&r&r&r\\
 2&*&rt&rht&rt&rt\\
 7&tI_2&t^{-1}&ht&t&t\\
 8&I_2&1&1&t^{-1}&h\\
 10&I_2&1&t^{-1}&h&1
 \end{array}
                                                               \tag{F.29.22}
\]
All entries use the original source indices;
in particular $10=\texttt{old\_-1\_4}$.

The five possible types are
$W,\partial,X_5,X_6,X_7$, with the classification
proved from the displayed maps above.
Source $0$ separates boundary from every
other type since $r_j\ne1$.
Two core types determine $(r_j,t_j)$ by
sources $0,7$; two boundary types do so by
sources $7,2$. For the same meridian type,
sources $0,7$ determine the same pair, with
the common $h$ cancelled for type $X_5$.
Every same-type comparison across distinct
factors therefore contradicts (F.29.9).

For the different types, core versus $X_7$
is excluded by source $8$, with scalars $1,h$.
Core versus $X_6$ is excluded by source $10$.
Core versus $X_5$ in factor $l$ would force
$t_l=1$ by source $10$. Inverse Frobenius then
gives all $t_j=1$, and source $7$ compares
$1$ with $h\ne1$.
Types $X_5,X_7$ are excluded by source $8$;
types $X_6,X_7$ by source $10$.
For type $X_5$ in factor $j$ versus $X_6$
in factor $l$, source $8$ forces $t_l=1$,
hence all $t_j=1$, and source $10$ then
compares $1$ with $h$.
This proves every cross-type comparison,
in either orientation. No extra condition
on $t-1$ or on the characteristic has entered
this chart.

\emph{Chart $c=1$, $r=1$: an essential line and actual seed evaluations.}
Here the character is
$(1,s,t,(st)^{-1},1,1,1)$, $h=st\ne1$,
$k_1=\mathbb F_p(h)$ and $k_0=k_1(t)=k_1(s)$.
The coefficient map is $r\mapsto1$ in
(F.29.19), sending $h$ to $st$ and retaining
the principal pivot $X_4$.
Since the field extension is proper, both
$s,t\notin k_1$. In particular
$s-1,t-1,t+1,h-1$ are nonzero, in every
characteristic, as are the units $s,t$.

Keep the six coordinates of the preceding
chart and put $b=(1,1,s,1,1,1)$.
The common kernel of the boundary and the
four character-one meridian covectors is
the line generated by $v_0=(0,1,-s^{-1},0,0,0)$.
The following basis and inverse are exact:
\[
 \begin{gathered}
 J=(v_0,e_2,e_5-e_2,e_4-e_2,e_1-e_2,e_6-e_2),\\
 J^{-1}(v)=(-sv_3,bv,v_5,v_4,v_1,v_6),\\
 J(z_0,z_\partial,z_6,z_5,z_1,z_7)
  =(z_1,z_0+z_\partial-z_6-z_5-z_1-z_7,
                                  -s^{-1}z_0,z_5,z_6,z_7).
 \end{gathered}                                                \tag{F.29.23}
\]
Substitution proves both composites to be the
identity. In these coordinates the actual
source matrices and one required lift identity are
\[
 \begin{gathered}
 T_l=\begin{pmatrix}\nu_l&a_l\\0&\Lambda_l\end{pmatrix},
 \qquad \Lambda_l\text{ diagonal},\\
 (\nu_0,\ldots,\nu_{10})=(1,1,t,1,s,1,1,t,1,s,1),\\
 \widetilde v_0=S_4v_0,\qquad
 A_2\widetilde v_0=t\widetilde v_0+st^2\delta.
 \end{gathered}                                                \tag{F.29.24}
\]
All eleven full matrices and the displayed
principal correction are independently verified.

For $j=1,2,5,7,9$, the additive commutators
$T_4T_j-T_jT_4$ kill the first line and have
image in that line. Their first rows on the
last five coordinates form $gH$, where
\[
 \begin{gathered}
 g=(s-1)(t-1),\qquad
 H=\begin{pmatrix}
 0&0&0&(st+1)(st+s-1)&-t^{-2}\\
 0&0&1&-(t+1)&0\\
 0&s/t&-t^{-2}&0&0\\
 (t+1)(st-t+1)/(st^3)&0&s&0&0\\
 -(s-1)/(st^2)&0&0&-(s-1)&0
 \end{pmatrix}.
 \end{gathered}                                                \tag{F.29.25}
\]
The independent signed-permutation expansion
of the determinant gives
\[
 \det H=\frac{(s-1)(t-1)(t+1)}{t^6},
 \qquad
 \det(gH)=\frac{(s-1)^6(t-1)^6(t+1)}{t^6}\ne0.
                                                               \tag{F.29.26}
\]
The factor proof preceding (F.29.23) justifies
every nonzero factor, including in characteristic
two where $t+1=t-1$.
Every nonzero subspace invariant under the
actual eleven sources contains
$L=k_0[v_0]$. Indeed, a vector with nonzero
last five coordinates is sent by some displayed
commutator to a nonzero first-line vector,
by (F.29.26). That vector remains in the
invariant subspace and spans $L$. A nonzero
subspace already contained in that line
equals $L$. This proves the essential-line
statement directly from the source operators.

Let $d_0$ have generator values $\widetilde v_0$.
For the actual negative seed
$a=u_{\texttt{cross\_0\_-1\_1}}$ of index $13$,
the word calculation and (F.29.24) give
\[
 \begin{gathered}
 \chi(a)=h^{-1},\qquad d_0(a)=h-1,\\
 \kappa=B_2(a)a^{-1}\in K,\qquad \chi(\kappa)=1,
                    \qquad d_0(\kappa)=1-h.
 \end{gathered}                                                \tag{F.29.27}
\]
In fact $d_0(B_2(a))=t(h-1)+ht(h^{-1}-1)=0$.
For a character-one word the principal
correction in (F.29.24) vanishes. Thus
\[
 \kappa_n=B_2^n(\kappa)\in K,\qquad
 \chi(\kappa_n)=1,\qquad
                       d_0(\kappa_n)=t^n(1-h)\quad(n\geq0).
                                                               \tag{F.29.28}
\]
The source preserves $K$, and the seed belongs
to the original eight negative seeds, so these
are actual elements of $K$. Independently, the
first two derivative rows
$D(a)(A_2-I_7)$ and $D(a)(A_2-I_7)A_2$ are
checked against (F.29.27)--(F.29.28), including
their vanishing on the principal vector.

\emph{The common restriction map and its exact kernel.}
Across all four charts consider the actual map
\[
 \Phi:\bigoplus_{j=0}^{m-1}Q_j
       \longrightarrow H^1_{\rm cont}(K,k_{0,\chi|K}),
       \qquad ([d_j])_j\longmapsto\sum_j[d_j|K].
                                                               \tag{F.29.29}
\]
It is well-defined because the characters
agree on $K$, and it commutes with every
source operator $[d]\mapsto[d\circ B_l]$.
Each individual restriction is injective
by F.24, since $\chi_j(R)=C\ne\{1\}$.

Choose a smallest-dimensional nonzero
source-invariant subspace $T\leq\ker\Phi$
under the supposition that the kernel is nonzero.
Each nonzero coordinate projection $p_j|T$
is injective, since its kernel is invariant.
Its image is a nonzero invariant subspace
with no proper nonzero invariant subspace.
Two nonzero projections therefore give the
exact mutually inverse intertwiners
\[
 p_l\circ(p_j|T)^{-1}:p_j(T)\longrightarrow p_l(T),
 \qquad p_j\circ(p_l|T)^{-1}:p_l(T)\longrightarrow p_j(T).
                                                               \tag{F.29.30}
\]
In the first three charts every pair of
possible types was excluded above for $j\ne l$.
There is therefore only one nonzero projection;
$T$ lies in a single factor, contradicting
that factor's F.24 restriction injectivity.
This proves $\Phi$ injective in those charts.

In the last chart the essential-line argument
puts $L_j$ inside each nonzero projection
image. Minimality makes that image equal to
$L_j$, so $T\leq\bigoplus_jL_j$.
Frobenius transport of (F.29.28) gives, on
these lines, the evaluation matrix at
$\kappa_0,\ldots,\kappa_{m-1}$ and its determinant
\[
 ((1-h)t_j^n)_{0\leq n,j<m},
 \qquad (1-h)^m\prod_{j<l}(t_l-t_j)\ne0.
                                                               \tag{F.29.31}
\]
All these evaluations descend to the common
target cohomology, since the words have
character one. The coefficient inverse is
explicit: write
\[
 \prod_{l\ne j}\frac{X-t_l}{t_j-t_l}
       =\sum_{n=0}^{m-1}h_{jn}X^n,
 \qquad z_j=(1-h)^{-1}\sum_{n=0}^{m-1}h_{jn}y_n.
                                                               \tag{F.29.32}
\]
It inverts $y_n=\sum_j(1-h)t_j^nz_j$;
evaluation of the displayed polynomial at
each $t_l$ proves the inverse composites.
Every denominator is nonzero by $h\ne1$
and the proved distinctness of the $t_j$.
Thus evaluation after $\Phi$ is injective
on $\bigoplus_jL_j$, contradicting the nonzero
$T$ in its kernel. This proves the injectivity
also in the last chart.

Zero summed restriction of cocycles now
forces each cohomology class to be zero,
so each cocycle is $b_j(\chi_j-1)$.
Its summed restriction is
$(\chi|K-1)\sum_jb_j$. Since $C\ne\{1\}$,
this vanishes exactly for $\sum_jb_j=0$.
The converse follows by direct restriction.
The principal vector in every factor is
nonzero, and the displayed parameter-sum
kernel has dimension $m-1$.
This proves (F.29.3), including the rank
$7m-(m-1)=6m+1$.

For the exact trace reduction of F.26, the
tuple $(d^{q_1^j})_j$ gives
\[
 \operatorname{Tr}_{k_0/k_1}(d|K)=0
 \quad\Longleftrightarrow\quad
 d=b(\chi-1)\ \hbox{with}\ \operatorname{Tr}_{k_0/k_1}(b)=0.
                                                               \tag{F.29.33}
\]
The forward implication is the proved cocycle
kernel, and the reverse is direct evaluation,
using $\chi(K)\subset k_1$.
F.26 sends every proper finite affine marginal
image through its explicit conjugation and
coefficient homomorphism to a nonprincipal
element of this trace kernel over the original
scalar field $k_0$. Those maps retain $\chi$,
so they retain (F.29.1), $C$, $k_1$ and $k_0$.
Equation (F.29.33) excludes that element.
This proves (F.29.2) for every remaining
negative representation, including those
with an original coefficient field larger
than $k_0$.

\emph{Integral evidence and the actual sign transport.}
The three symbolic producers retain all full
source matrices and basis maps for the generic
four-parameter chart and the two pivot-$X_4$
charts. Their exact check counts are $754$,
$691$ and $956$.
The separately structured three-chart verifier
imports neither a producer nor an algebra
package. It reconstructs all eleven old-source
and all eight negative-seed braid words in
each of those three integral Laurent rings.
Its $3246$ checks include (F.29.7), both basis
inverses, every invariant inclusion and diagonal
quotient, the comparison tables, the integral
annihilators, the five printed commutator rows,
their $120$-term determinant expansion, the
source-$2$ principal correction and the actual
seed evaluations.
The additional $r=1,c\ne1$ verifier uses the
same reviewed integer Laurent arithmetic,
reconstructs all eleven original words again,
and records $996$ checks of the new basis,
every quotient character and (F.29.16).
Root's separate exceptional audit reconstructs
sources $1,2,7$ directly and checks the exact
coefficient map and both residual identities
in (F.29.18), with $20$ additional recorded
certificate assertions.

These are integral coefficient identities.
The only nonmonomial principal localizations
are $1-c$ and $1-h$ in their respective
charts; their nonvanishing was proved before
each quotient was used. All other determinant
and interpolation factors have their explicit
nonzero proofs above. The exceptional constant
$2$ is inverted only after (F.29.17) proves
odd characteristic. The evidence manifest
binds the source words, programs, outputs,
coefficient inputs and completed dependencies
to this exact locus and statement.

Finally, a positive representation precomposes
with $\mathcal B_+:V_-\to V_+$ to have precisely
the negative character just proved in the
basis $Y_i^+=\mathcal B_+(X_i)$.
This continuous isomorphism and its inverse
carry $\mathcal K_-,R_-$ onto their positive
counterparts by F.23.13. Equality of the two
negative images therefore gives equality of
the two original positive images through
those exact maps. This proves (F.29.2) for
both signs.
\end{proof}
